\subsection{NOvA Dual-Baseline Search for Active-to-Sterile Neutrino Oscillations using Neutrino- and Antineutrino-Enriched Samples --- arXiv:2609.08900}

\textbf{Authors:} NOvA Collaboration

\paragraph{主な主張}
NOvA の near/far detector を用いて neutrino-mode と antineutrino-mode のデータを同時解析した結果、$3+1$ sterile-neutrino model の有意な兆候は得られず、従来実験で許されていた領域の一部、特に IceCube が許していた領域の大部分を排除した \cite{NOvA:2026sterile}。

\paragraph{新規性}
NOvA として初めて neutrino/antineutrino の両 beam mode を含む dual-baseline fit を行い、Near Detector での eV-scale oscillation distortion と Far Detector での長基線情報を同時に用いて sterile-neutrino 感度を拡張した点が新しい。

\paragraph{位置付け}
eV-scale sterile sector に対する加速器ベースの強い制限を与える一方、単一 sterile state を仮定する $3+1$ 解析であるため KK tower への適用には専用の recast が必要であり、TeV 領域の Earth-matter resonance を利用する IceCube と相補的な probe と位置付けられる。
